\chapter{Firmware and Sensor Signal Processing}

The firmware exists in two roles, sharing a common signal-processing core: the
\emph{slave} sketch (thumb, index, ring, pinky) and the \emph{master} sketch
(middle finger). Both are written for the Arduino-ESP32 framework.

\section{Per-Finger Signal Processing}

Every module, master or slave, processes its own IMU identically.

\subsection{Acquisition and Filtering}

The MPU6050 is read over I\textsuperscript{2}C at 25\,Hz. The on-sensor digital
low-pass filter is set to 10\,Hz, which removes high-frequency mechanical noise
before sampling. A bus-recovery routine bit-bangs the clock line and
re-initialises the sensor if a sequence of reads fails, so a transient
I\textsuperscript{2}C stall does not permanently disable a finger.

\subsection{Orientation Estimation}

Pitch and roll are estimated by fusing the accelerometer and gyroscope with a
two-state Kalman filter \cite{kalman} per axis, with the angle and the gyro
bias as states. When the magnitude of the measured acceleration departs from
$1\,g$ (indicating linear motion rather than pure gravity), the filter trusts
the gyroscope more and gates out the unreliable accelerometer tilt. Yaw, which
a six-axis IMU cannot reference absolutely, is obtained by integrating the
bias-corrected $z$-axis angular rate with a small dead-band, and is reset to
zero after a sustained period of stillness to bound its drift.

\subsection{Calibration and Bias Tracking}

At start-up each module performs a stillness calibration: it averages several
hundred samples to establish the zero references for pitch and roll and the
initial gyroscope bias, rejecting the attempt if the measured variance shows
the hand was moving. During operation, whenever the module detects stillness
(low angular rate and acceleration close to $1\,g$), it slowly updates the
gyroscope bias with an exponential moving average. This continuously tracks the
bias drift caused by the microcontroller warming the adjacent sensor, which is
significant for a low-cost gyroscope and is one reason motion-based features
later proved more robust than absolute-orientation features.

\subsection{Packet Format}

Each module produces, per cycle, a packet of nine fields: the three fused
angles (pitch, roll, yaw), the three bias-corrected angular rates
($g_x,g_y,g_z$), and the three accelerations ($a_x,a_y,a_z$). A slave sends
this as a small fixed-size structure tagged with its finger identifier and a
sequence number.

\section{The Slave Sketch}

A slave runs a deadline-scheduled 25\,Hz loop. Each cycle it reads and fuses
its IMU, builds its packet, and transmits it by ESP-NOW \emph{unicast} to the
master's address on a fixed Wi-Fi channel. A staggered start-up delay
per finger reduces the chance that all four slaves transmit in the same
instant. The module also prints a periodic diagnostic line (read success rate,
transmit success/failure counts, current angles) over USB for bench debugging.

\section{The Master Sketch}

The master is dual-purpose, selected at compile time:

\begin{itemize}
    \item \textbf{CSV mode} streams every received packet, plus its own, as
    comma-separated lines over USB at 921600\,baud. This feeds the data-%
    collection script (Chapter~5). The high baud rate is necessary: at
    115200\,baud the five-finger stream saturates the USB serial buffer, the
    blocking print inside the receive callback backs up, and 15--30\% of
    packets are dropped—a failure that was diagnosed and fixed by raising the
    baud rate.
    \item \textbf{BLE mode} exposes a GATT service with a notify
    characteristic. Each 25\,Hz cycle the master reads its own IMU, packs the
    five fingers' nine fields into a 180-byte payload (45 little-endian 32-bit
    floats), and sends it as a BLE notification to the connected phone. A
    finger whose data has not arrived recently is sent as zeros so the payload
    is always well-formed.
\end{itemize}

\subsection{Aggregation and Channel Order}

The master keeps the latest packet from each finger in a lock-protected table.
The 45 channels are emitted in a fixed finger-major order—thumb, index, middle,
ring, pinky—each contributing $P,R,Y,g_x,g_y,g_z,a_x,a_y,a_z$. This exact order
is mirrored in the data-collection script and in the phone application, so that
the model sees identically ordered channels everywhere.

\subsection{Radio Coexistence and the BLE MTU}

The master runs ESP-NOW (Wi-Fi physical layer) reception and BLE transmission
on a single radio, which the ESP32-C3 time-slices. At 25\,Hz this is
comfortable. One non-obvious requirement is the BLE \emph{maximum transmission
unit}: the default 23-byte MTU carries only 20 payload bytes per notification,
far short of the 180-byte feature vector, so the phone must negotiate a larger
MTU (247 bytes requested, 256 granted) immediately after connecting. Without
this negotiation the notifications are truncated and silently rejected by the
phone—a failure that initially presented as ``no data'' despite a healthy
radio link.
